\subsection{AUTOMORPHISMS OF A SPLITTABLE SEMI-SIMPLE LIE ALGEBRA}

PROPOSITION 5. Assume that g is splittable. The group \(\operatorname{Aut}_{0}(\mathfrak{g})\) operates simply-transitively on the set of framings of g.

Let \(e_1 = (\mathfrak{g}, \mathfrak{h}_1, \mathrm{B}_1, (X_\alpha^1)_{\alpha \in \mathrm{B}_1})\) , \(e_2 = (\mathfrak{g}, \mathfrak{h}_2, \mathrm{B}_2, (X_\alpha^2)_{\alpha \in \mathrm{B}_2})\) be two framings of \(\mathfrak{g}\) . There exists at least one element of \(\operatorname{Aut}_0(\mathfrak{g})\) that transforms \(e_1\) into \(e_2\) (Prop. 1 and Prop. 4). Let \(\bar{k}\) be an algebraic closure of \(k\) . There exists an element of \(\operatorname{Aut}_e(\mathfrak{g} \otimes_k \bar{k})\) that transforms \(\mathfrak{h}_1 \otimes_k \bar{k}\) into \(\mathfrak{h}_2 \otimes_k \bar{k}\) (Chap. VII, §3, no. 2, Th. 1). Hence, by Prop. 4 and Cor. 2 of Prop. 2, there exists an element \(\varphi\) of \(\operatorname{Aut}_e(\mathfrak{g} \otimes_k \bar{k})\) that transforms the framing \((\mathfrak{g} \otimes_k \bar{k}, \mathfrak{h}_1 \otimes_k \bar{k}, \mathrm{B}_1, (X_\alpha^1)_{\alpha \in \mathrm{B}_1})\) of \(\mathfrak{g} \otimes_k \bar{k}\) into the framing \((\mathfrak{g} \otimes_k \bar{k}, \mathfrak{h}_2 \otimes_k \bar{k}, \mathrm{B}_2, (X_\alpha^2)_{\alpha \in \mathrm{B}_2})\) . Since \(\mathfrak{h}_1\) and the \(X_\alpha^1\) (resp. \(\mathfrak{h}_2\) and the \(X_\alpha^2\) ) generate \(\mathfrak{g}_1\) (resp. \(\mathfrak{g}_2\) ), we have \(\varphi(\mathfrak{g}_1) = \mathfrak{g}_2\) , so \(\varphi\) is of the form \(\psi \otimes 1\) where \(\psi \in \operatorname{Aut}_0(\mathfrak{g})\) , and \(\psi\) transforms \(e_1\) into \(e_2\) .

COROLLARY 1. Let \((\mathfrak{g},\mathfrak{h},\mathrm{B},(X_{\alpha})_{\alpha\in\mathrm{B}})\) be a framing of g, and G the group (isomorphic to \(\operatorname{Aut}(\mathrm{R},\mathrm{B})\) ) of automorphisms of g that leave this framing invariant. Then \(\operatorname{Aut}(\mathfrak{g})\) is the semi-direct product of G and \(\operatorname{Aut}_{0}(\mathfrak{g})\) .

Indeed, every element of \(\operatorname{Aut}(\mathfrak{g})\) transforms \((\mathfrak{g},\mathfrak{h},\mathrm{B},(X_{\alpha})_{\alpha\in\mathrm{B}})\) into a framing of g. By Prop. 5, every coset of \(\operatorname{Aut}(\mathfrak{g})\) modulo \(\operatorname{Aut}_{0}(\mathfrak{g})\) meets G in exactly one point. Q.E.D.

It follows from Cor. 1 that the group \(\operatorname{Aut}(\mathfrak{g})/\operatorname{Aut}_{0}(\mathfrak{g})\) can be identified with \(\operatorname{Aut}(\mathrm{R},\mathrm{B})\) , and is isomorphic to the group of automorphisms of the Dynkin graph of R.

COROLLARY 2. \(\operatorname{Aut}(\mathfrak{g}) = \operatorname{Aut}_0(\mathfrak{g})\) when \(\mathfrak{g}\) is a splittable simple Lie algebra of type \(\mathrm{A}_1\) , \(\mathrm{B}_n\) ( \(n \geq 2\) ), \(\mathrm{C}_n\) ( \(n \geq 2\) ), \(\mathrm{E}_7\) , \(\mathrm{E}_8\) , \(\mathrm{F}_4\) , \(\mathrm{G}_2\) . The quotient \(\operatorname{Aut}(\mathfrak{g}) / \operatorname{Aut}_0(\mathfrak{g})\) is of order 2 when \(\mathfrak{g}\) is of type \(\mathrm{A}_n\) ( \(n \geq 2\) ), \(\mathrm{D}_n\) ( \(n \geq 5\) ), \(\mathrm{E}_6\) ; it is isomorphic to \(\mathfrak{S}_3\) when \(\mathfrak{g}\) is of type \(\mathrm{D}_4\) .

This follows from Cor. 1 and Chap. VI, Plates I to IX.

Remarks. 1) Let \(e_1 = (\mathfrak{g}, \mathfrak{h}_1, \mathrm{B}_1, (X_\alpha^1)_{\alpha \in \mathrm{B}_1})\) , \(e_2 = (\mathfrak{g}, \mathfrak{h}_2, \mathrm{B}_2, (X_\alpha^2)_{\alpha \in \mathrm{B}_2})\) , \(e_2' = (\mathfrak{g}, \mathfrak{h}_2, \mathrm{B}_2, (Y_\alpha^2)_{\alpha \in \mathrm{B}_2})\) be framings of \(\mathfrak{g}\) , and \(s\) (resp. \(s'\) ) an element of \(\operatorname{Aut}_0(\mathfrak{g})\) that transforms \(e_1\) to \(e_2\) (resp. \(e_2'\) ). Then \(s|_{\mathfrak{h}_1} = s'|_{\mathfrak{h}_1}\) . Indeed, \(s'^{-1}s \in \operatorname{Aut}_0(\mathfrak{g}, \mathfrak{h}_1)\) and \(s'^{-1}s(\mathrm{B}_1) = \mathrm{B}_1\) , so \(\varepsilon(s'^{-1}s) = 1\) .

\begin{enumerate}
\def\labelenumi{\arabic{enumi})}
\setcounter{enumi}{1}
\tightlist
\item
  Let X be the set of pairs \((\mathfrak{h},\mathrm{B})\) where h is a splitting Cartan subalgebra of g and B a basis of the root system of \((\mathfrak{g},\mathfrak{h})\) . If \(x=(\mathfrak{h},\mathrm{B})\) and \(x'=(\mathfrak{h}',\mathrm{B}')\) are two elements of X, there exists \(s\in\operatorname{Aut}_{0}(\mathfrak{g})\) that transforms x into \(x'\) (Prop. 5), and the restriction \(s_{x',x}\) of s to h does not depend on the choice of s (Remark 1). In particular, \(s_{x'',x'}\circ s_{x',x}=s_{x'',x}\) if \(x,x',x''\in X\) , and \(s_{x,x}=1\) . The set of families \((h_{x})_{x\in\mathrm{X}}\) satisfying the conditions
\end{enumerate}

\[
a) h _ {x} \in \mathfrak {h} \text {   if   } x = (\mathfrak {h}, B)
\]

\[
b) s _ {x ^ {\prime}, x} (h _ {x}) = h _ {x ^ {\prime}} \text { if } x, x ^ {\prime} \in X
\]

is in a natural way a vector space \(\mathfrak{h}(\mathfrak{g})\) which we sometimes call the canonical Cartan subalgebra of g. For \(x = (\mathfrak{h}, \mathrm{B})\) and \(x' = (\mathfrak{h}', \mathrm{B}')\) , \(s_{x', x}\) takes B to \(B'\) , and hence the root system of \((\mathfrak{g}, \mathfrak{h})\) to that of \((\mathfrak{g}, \mathfrak{h}')\) ; it follows that the dual \(\mathfrak{h}(\mathfrak{g})^*\) of \(\mathfrak{h}(\mathfrak{g})\) is naturally equipped with a root system \(\mathrm{R}(\mathfrak{g})\) and with a basis \(\mathrm{B}(\mathfrak{g})\) of \(\mathrm{R}(\mathfrak{g})\) . We sometimes say that \(\mathrm{R}(\mathfrak{g})\) is the canonical root system of g and that \(\mathrm{B}(\mathfrak{g})\) is its canonical basis. The group \(\mathrm{Aut}(\mathfrak{g})\) operates on \(\mathfrak{h}(\mathfrak{g})\) leaving \(\mathrm{R}(\mathfrak{g})\) and \(\mathrm{B}(\mathfrak{g})\) stable; the elements of \(\mathrm{Aut}(\mathfrak{g})\) that operate trivially on \(\mathfrak{h}(\mathfrak{g})\) are those of \(\mathrm{Aut}_0(\mathfrak{g})\) .

PROPOSITION 6. Let \(\mathfrak{h}\) be a splitting Cartan subalgebra of \(\mathfrak{g}\) . We have, with the notations in no. 1, \(\mathrm{Aut}_0(\mathfrak{g}) = \mathrm{Aut}_e(\mathfrak{g}).\mathrm{Ker}\varepsilon = \mathrm{Aut}_e(\mathfrak{g}).f(\mathrm{T}_Q)\) .

By §3, no. 3, Cor. of Prop. 10, \(\mathrm{Aut}_0(\mathfrak{g}) = \mathrm{Aut}_e(\mathfrak{g}).\mathrm{Aut}_0(\mathfrak{g},\mathfrak{h})\) . On the other hand, \(\varepsilon (\mathrm{Aut}_e(\mathfrak{g},\mathfrak{h}))\supset \mathrm{W}(\mathrm{R})\) by §2, no. 2, Cor. of Th. 2, so \(\mathrm{Aut}_0(\mathfrak{g},\mathfrak{h}) = \mathrm{Aut}_e(\mathfrak{g},\mathfrak{h}).\mathrm{Ker}\varepsilon\) .

Remark 3. Prop. 6 shows that the canonical homomorphism

\[
\iota : \mathrm{T} _ {\mathrm{Q}} / \operatorname{Im} (\mathrm{T} _ {\mathrm{P}}) \to \operatorname{Aut} _ {0} (\mathfrak {g}) / \operatorname{Aut} _ {e} (\mathfrak {g}),
\]

induced by the diagram in no. 2, is surjective. In particular, \(\mathrm{Aut}_e(\mathfrak{g})\) contains the derived group of \(\mathrm{Aut}_0(\mathfrak{g})\) ; we shall see (§11, no. 2, Prop. 3) that they are actually equal. Moreover, it can be shown that \(\iota\) is injective, in other words that

\[
f (\mathrm{T} _ {\mathrm{Q}}) \cap \operatorname{Aut} _ {e} (\mathfrak {g}) = f ^ {\prime} (\mathrm{T} _ {\mathrm{P}}),
\]

(cf.~§7, Exerc. 26 d)).

PROPOSITION 7. Let g be a splittable semi-simple Lie algebra, b a Borel subalgebra of g, and \(p_{1}\) and \(p_{2}\) distinct parabolic subalgebras of g containing b. Then \(p_{1}\) and \(p_{2}\) are not conjugate under \(\operatorname{Aut}_{0}(\mathfrak{g})\) .

We can assume that \(k\) is algebraically closed. Let \(s \in \mathrm{Aut}_0(\mathfrak{g})\) be such that \(s(\mathfrak{p}_1) = \mathfrak{p}_2\) . Let \(\mathfrak{h}\) be a Cartan subalgebra of \(\mathfrak{g}\) contained in \(\mathfrak{b} \cap s(\mathfrak{b})\) (§3, no. 3, Prop. 10). Since \(\mathfrak{h}\) and \(s(\mathfrak{h})\) are Cartan subalgebras of \(s(\mathfrak{b})\) , there exists \(u \in [\mathfrak{b}, \mathfrak{b}]\) such that \(e^{\mathrm{ad}u}(\mathfrak{h}) = s(\mathfrak{h})\) (Chap. VII, §3, no. 4, Th. 3). Replacing \(s\) by \(e^{-\mathrm{ad}u}s\) , we are reduced to the case in which \(s(\mathfrak{h}) = \mathfrak{h}\) , and \(s\) then induces on \(\mathfrak{h}\) an element \(\sigma\) of the Weyl group W of \((\mathfrak{g}, \mathfrak{h})\) (Prop. 4). Let C be the Weyl chamber corresponding to \(\mathfrak{b}\) . Then \(\mathfrak{p}_1\) and \(\mathfrak{p}_2\) correspond to facets \(F_1\) and \(F_2\) of \(\mathfrak{h}_{\mathbf{R}}\) contained in the closure of C. We have \(\sigma(F_1) = F_2\) . Since \(\sigma \in W\) , this implies that \(F_1 = F_2\) (Chap. V, §3, no. 3, Th. 2) so \(\mathfrak{p}_1 = \mathfrak{p}_2\) .

Remark 4. Let g be a splittable semi-simple Lie algebra, P the set of parabolic subalgebras of g, a set on which \(\operatorname{Aut}_{0}(\mathfrak{g})\) operates. Retain the notations of Remark 2. Let \(\Sigma\) be a subset of \(\mathrm{B}(\mathfrak{g})\) . Giving \(\Sigma\) is equivalent to giving, for every \(x = (\mathfrak{h}, \mathrm{B}) \in \mathrm{X}\) , a subset \(\Sigma_{x}\) of B, such that \(s_{x',x}\) takes \(\Sigma_{x}\) to \(\Sigma_{x'}\) for any \(x, x' \in X\) . Let \(p_{x}\) be the parabolic subalgebra of g corresponding to \(\Sigma_{x}\) (§3, no. 4, Remark). The orbit of \(p_{x}\) under \(\operatorname{Aut}_{0}(\mathfrak{g})\) is the set of \(p_{x'}\) for \(x' \in X\) . This defines a map from \(\mathfrak{P}(\mathrm{B}(\mathfrak{g}))\) to \(\mathcal{P}/\operatorname{Aut}_{0}(\mathfrak{g})\) . This map is surjective by the Remark of §3, no. 4, and injective by Prop. 7.

